\begin{Exercise}[label={ex:facebook}]
  \Question Si definiscano le seguenti variabili che specificano le strutture dati per rappresentare i messaggi scambiati attraverso Facebook e ai profili degli utenti:
  		\begin{itemize}
  				\item una struttura che definisca un'utente, con nome, cognome e e-mail, tutti composti da caratteri alfanumerici;
  				\item una struttura che definisca un messaggio, con un contenuto alfanumerico, un mittente scelto tra gli utenti, un numero fissato di destinatari e i destinatari, anch'essi scelti tra gli utenti (fino ad un massimo di 256).
  		\end{itemize}

		\Question Definire, usando il linguaggio C, un'appropriata variabile per memorizzare le informazioni relative a 10 utenti e 10 messaggi. Queste variabili dovranno essere chiamate \texttt{utenti\_facebook} e \texttt{messaggi\_inviati}.

		\Question Scrivere in linguaggio C, aggiungendo eventualmente opportune dichiarazioni di variabili, un frammento di codice che permetta di visualizzare a video l'indirizzo email di tutti gli utenti che hanno spedito almeno un messaggio e che non siano tra i destinatari di tale messaggio. Si presupponga che la variabile \texttt{messaggi\_inviati} sia inizializzata correttamente con le informazioni relative a 10 messaggi. Si presupponga inoltre che l'indirizzo email sia identificatore univoco di un utente: pertanto si può utilizzare l'indirizzo email per verificare se due utenti sono lo stesso utente.
		
\end{Exercise}

\begin{Answer}[ref={ex:facebook}]

\begin{clangsmall}
typedef struct {
  char nome[24];
  char cognome[24];
  char email[64];
} Utente;

typedef struct {
  char contenuto[256];
  Utente mittente;
  int numDestinatari;
  Utente destinatari[256];
} Messaggio;
\end{clangsmall}

\begin{clangsmall}
Utente utenti_facebook[10];
Messaggio messaggi_inviati[10];		
\end{clangsmall}

Gli utenti che hanno spedito almeno un messaggio sono necessariamente presenti come mittenti nella variabile \texttt{messaggi\_inviati}:

\begin{clangsmall}
void main () {
  int i, j;
  Utente utenti_facebook[10];
  Messaggio messaggi_inviati[10];		

  for (i = 0; i < 10; i++) {
    j = 0;

    //ricerca destinatario.email != mittente.email con strcmp()
    while (j < messaggi_inviati[i].numDestinatari &&
              strcmp(messaggi_inviati[i].mittente.email,
                     messaggi_inviati[i].destinatari[j].email) != 0) {
      j++;
    }

    if (j == messaggi_inviati[i].numDestinatari) //non trovato
      printf("%s", Messaggi[i].mittente.email);
}
    \end{clangsmall}
	
\end{Answer}